\section{Q5 — Raft Failure Test Cases}

To evaluate fault tolerance in our Raft implementation, we performed five
failure-oriented tests while monitoring the live cluster logs. We used
\texttt{docker compose stop/start} to control individual nodes and filtered
logs using \texttt{grep} to focus on election, leader, and commit events.
These tests focus on node departure, node return, leader failure, leader
recovery, and the quorum rule for a 5-node cluster.

\textbf{Commands used:}
\begin{verbatim}
docker compose logs -f api-server-a api-server-b api-server-c \
  api-server-d api-server-e \
  | grep -E "RAFT|CANDIDATE|LEADER|RequestVote|lost majority|\
            isolated|failed|unreachable|COMMIT|APPLY|Stopping|\
            Restarting"

docker compose stop <node-name>
docker compose start <node-name>
\end{verbatim}

% ── Test Case 1 ──────────────────────────────────────────
\subsection*{Test Case 1 — Non-Leader Node Exits}

\textbf{Setup:} Stop one follower using \texttt{docker compose stop api-server-<X>}.

\textbf{Expected behavior:} Because 4 of the 5 nodes remain alive, the leader
keeps control of the cluster and Raft continues to make progress. No re-election
is required because the failed node is not the leader.

\begin{figure}[H]
  \centering
  \includegraphics[width=0.95\textwidth]{images/followerFail.png}
  \caption{Test 1: A non-leader node exits the cluster while the remaining nodes continue operating normally.}
  \label{fig:q5_t1}
\end{figure}

% ── Test Case 2 ──────────────────────────────────────────
\subsection*{Test Case 2 — Non-Leader Node Rejoins}

\textbf{Setup:} Restart the stopped follower using \texttt{docker compose start api-server-<X>}.

\textbf{Expected behavior:} The returning node resumes participation as a
follower, receives heartbeats / log updates from the leader, and re-enters the
cluster without forcing a new election.

\begin{figure}[H]
  \centering
  \includegraphics[width=0.95\textwidth]{images/followerReturn.png}
  \caption{Test 2: A previously stopped follower rejoins the Raft cluster and resumes normal operation.}
  \label{fig:q5_t2}
\end{figure}

% ── Test Case 3 ──────────────────────────────────────────
\subsection*{Test Case 3 — Leader Failure}

\textbf{Setup:} Identify the current leader from the logs, then stop it using
\texttt{docker compose stop <leader-node>}.

\textbf{Expected behavior:} Once heartbeats from the leader disappear, one of
the remaining nodes times out, becomes a candidate, sends \texttt{RequestVote}
RPCs, and is elected as the new leader by a majority of the cluster.

\begin{figure}[H]
  \centering
  \includegraphics[width=0.95\textwidth]{images/leaderFail.png}
  \caption{Test 3: When the leader fails, the remaining nodes hold a new election and choose a replacement leader.}
  \label{fig:q5_t3}
\end{figure}

% ── Test Case 4 ──────────────────────────────────────────
\subsection*{Test Case 4 — Old Leader Rejoins After Failure}

\textbf{Setup:} Restart the previously failed leader using
\texttt{docker compose start <leader-node>}.

\textbf{Expected behavior:} The old leader does not immediately reclaim
leadership. Instead, it rejoins as a follower and begins receiving heartbeats
from the node that won the new election.

\begin{figure}[H]
  \centering
  \includegraphics[width=0.95\textwidth]{images/leaderReturn.png}
  \caption{Test 4: The failed leader rejoins the cluster and returns as a follower under the new leader.}
  \label{fig:q5_t4}
\end{figure}

% ── Test Case 5 ──────────────────────────────────────────
\subsection*{Test Case 5 — Quorum Rule Demonstration (2 Failures vs. 3 Failures)}

\textbf{Setup:} Stop two follower nodes using \texttt{docker compose stop}, observe
the cluster continuing to operate, then stop a third node.

\textbf{Expected behavior:}
\begin{itemize}
  \item With \textbf{2 failed nodes}, the cluster still has \textbf{3/5} nodes alive, which is a majority, so Raft can still commit operations.
  \item With \textbf{3 failed nodes}, only \textbf{2/5} nodes remain, which is not a majority, so the cluster loses quorum and can no longer safely make progress.
\end{itemize}

This test directly demonstrates the \(2k+1\) property of a 5-node Raft system:
it tolerates up to 2 failures, but not 3.

\begin{figure}[H]
  \centering
  \includegraphics[width=0.48\textwidth]{images/twoFails.png}\hfill
  \includegraphics[width=0.48\textwidth]{images/threeFails.png}
  \caption{Test 5: Quorum rule demonstration. Left: with two failed nodes, the remaining 3/5 nodes still form a majority and Raft continues. Right: with three failed nodes, quorum is lost and the cluster stalls.}
  \label{fig:q5_t5}
\end{figure}
